\section{Mechanistic Characterization: Testing the Causal Chain}
\label{sec:mechanism}

A design chain is only as credible as the measurements that can falsify each
link. This section reviews the toolbox in the order the chain itself
proposes -- factor isolation first, then carrier dynamics, then atomic-scale
and operando interface probes -- and adopts calibrated language throughout:
spectra \emph{show} observables, global analysis \emph{assigns} components
under a model, and matched activity correlations \emph{are consistent with}
rather than prove causality~\cite{choi2024understanding}.

\subsection{Factor deconvolution before spectroscopy}

Because COF photocatalysts co-vary crystallinity, porosity, absorption,
wettability, and charge extraction, single-material reports rarely identify
which factor moved. The structure--property--activity protocol of
Ghosh and co-workers -- quasi-independent variation of factors across the
S-COF/FS-COF benchmark family -- identified crystallinity and charge
extraction as the prime HER determinants and remains the template for
benchmarking claims~\cite{ghosh2020identification,wang2018sulfone}.
Substituent series on fixed backbones extend the same logic: six-membered
substituent sets isolate electronic effects from framework geometry
~\cite{jiang2024significant}, and matched isoreticular linker series
separate linker length from core identity~\cite{dong2025synergistic}.
Reviews of the crystallinity--stability dilemma emphasize that order
parameters must be reported alongside rates for any comparison to be
meaningful~\cite{ran2024progress,xiao2024molecular}.

\subsection{Carrier dynamics: from femtoseconds to seconds}

Time-resolved spectroscopy provides the quantitative observables of the
middle links. Femtosecond transient absorption (fs-TA) resolves exciton
formation, dissociation, and trapping; in S-scheme heterojunctions it
directly tracks interfacial electron flow, as shown for ZnSe quantum
dot/COF~\cite{zhao2024investigating} and FS-COF/WO$_3$ hybrids
~\cite{zhang2025revealing}. Time-resolved photoluminescence, surface
photovoltage, and wavelength-dependent photocurrent action spectra together
assigned the enhanced activity of g-C$_3$N$_4$/Ni-CAT junctions to transfer
of high-energy-level electrons rather than to absorption gains
~\cite{chu2024controlled}. Nanoscale COF colloids show molecule-like
excitonic behavior in photoluminescence and transient absorption, tying
carrier physics to particle dimension~\cite{zhao2024nanoscale}, and
lifetime engineering through D--A pairing is read out as dramatically
prolonged carrier lifetimes accompanying the 82.6\% AQE benchmark
~\cite{li2022covalent}. The historical caution from oxide photocatalysis
applies unchanged: lifetime changes alone are not productive charge
transfer unless they correlate with quantum yield under matched intensity
~\cite{tang2008mechanism,ma2022charge}.

\subsection{Atomic-scale and operando interface probes}

Where the chain meets the cocatalyst, aberration-corrected HAADF-STEM and
X-ray absorption fine structure establish speciation: atomically dispersed
Pt on TpPa-1 was shown to bind as a six-coordinated C$_3$N--Pt--Cl$_2$
species, with DFT connecting that geometry to the activity gain
~\cite{dong2021platinum}. Soft X-ray absorption identified
halogen-adjacent carbons as the active H$^*$ binding sites in Cl-decorated
frameworks~\cite{ma2025engineering}, and in situ XPS traced the S-scheme
electron pathway in fully conjugated COF/C$_3$N$_4$ junctions
~\cite{liu2025in}. Mechanistic assignment of the Volmer--Tafel route with a
Tafel rate-determining step on Ru single-atom/nanocluster hybrids shows how
kinetic analysis complements spectroscopy~\cite{qu2024covalent}. FT-IR and
$^1$H~NMR located hydrogen-bonded nanoreactor encapsulation inside COF
pores~\cite{pan2025reducing}, while ESIPT-induced photoisomerization of
keto-enamine linkages -- invisible to steady-state diffraction -- was caught
by excited-state spectroscopy, reshaping how activity differences within
the Tp family should be interpreted~\cite{weng2023effect}.

\subsection{What counts as mechanistic evidence}

Three practices from the strongest studies deserve adoption as field norms.
First, matched controls: physical blends and amorphous analogues distinguish
framework effects from composition effects~\cite{xu20202d,mi2021covalent}.
Second, multi-probe convergence: charge-separation claims supported jointly
by fs-TA kinetics, photoelectrochemical response, and activity correlations
~\cite{shen2023efficient,he2023integrated,zhang2023stable} are categorically
stronger than single-technique assignments. Third, protocol-anchored
quantum yields: AQE at stated wavelengths -- 6.7\% at 420~nm for alkene
linkages~\cite{mo2020alkene}, 2.53\% for sp$^2$c overall splitting
~\cite{cheng2023fully}, 9.39\% under piezo-assisted operation
~\cite{gao2025piezoceramics} -- remain the only observables portable across
laboratories, a theme Section~\ref{sec:performance} develops
quantitatively~\cite{zhong2023oligo}.
